% !TeX root = ../main.tex

\xchapter{面向长链条平面几何推理的两阶段增强方法}{A Two-Stage Enhancement Method for Long-Chain Reasoning in Plane Geometry}

为了解决现有模型在长链条推理中面临的性能瓶颈，本章提出了一种两阶段增强方法。该方法首先在纯文本环境中训练模型的基础推理能力，减少视觉信息干扰；随后，在图文联合环境中加入符号规则的步骤级反馈机制，从而提高推理稳定性，缓解误差累积与推理失稳问题。不同数据集上的实验结果验证了该方法的有效性。

% ==========================================================
\xsection{引言}{Introduction}

在 GeoLStep 基准上的基线实验表明，现有几何求解模型和大语言模型在较短推理链条上通常能够取得较好的效果，但当题目需要经过更多中间结论、依赖更复杂的条件关系时，模型性能会明显下降。这一现象说明，长链条平面几何推理的困难并不只在于答案求解本身，更在于模型能否在图文联合条件下持续维护正确的几何关系，并保持推理过程的连贯性与稳定性。

跨模态长链条推理中的逻辑失稳与误差累积风险，主要体现在两个方面。首先，平面几何题目包含文本条件和图形结构，图文信息的表达差异会影响模型对跨模态信息的统一建模~\cite{laurenccon2024matters}；随着推理链条延长，视觉信息的利用可能减弱，从而导致推理偏移~\cite{sun2025mitigating}。其次，长链条生成过程缺乏有效的中间状态约束，推理中的小错误可能会随着推理链的展开逐渐放大，影响最终结果。已有研究表明，中间推理噪声削弱了链式推理效果~\cite{zhou2024can}，而引入解释链等过程信息有助于提升多模态推理表现~\cite{lu2022learntoexplain}。

这些挑战表明，直接在图文联合输入下进行统一训练并不理想。模型需要同时处理图像识别、文本理解和多步推导等多种要求，训练难度较高。对于长链条几何题而言，模型首先需要具备稳定的条件整理和连续推导能力；而图像信息的引入虽然能够补充题目条件，但也会增加跨模态理解的复杂性。当二者在同一阶段中共同优化时，模型容易受到视觉信息噪声和图文关系不确定性的影响，导致推理能力的学习不够充分。因此，本章采用两阶段训练策略：先在纯文本环境中强化模型的基础推理能力，再在图文联合环境中加入过程级反馈，对多模态推理过程进行进一步约束。

为解决上述问题，本章提出了 GeoTSR（Geometry Two-Stage Reinforcement），一种面向长链条平面几何推理的两阶段增强方法。该方法在第一阶段将几何题目转化为纯文本形式，减少图像输入对训练过程的干扰，帮助模型集中学习条件整理、关系推导等推理能力；第二阶段则重新引入图文联合输入，并通过基于符号规则的步骤级反馈机制，对生成的中间结论进行检验与约束，从而提升推理能力，减少无效推理和路径偏移。实验结果表明，GeoTSR 在 MathVerse 和 GeoLStep 数据集上的表现优于基线方法，验证了其在长链条几何推理任务中的有效性。

% ==========================================================
\xsection{总体框架}{Overall Framework}

GeoTSR 的整体框架如图~\ref{fig:c4_framework} 所示，主要包括纯文本环境下的长链条推理增强、多模态条件下的过程级强化，以及贯穿训练过程的弱符号验证器三个部分。整体上，GeoTSR 采用由单模态推理能力建模到多模态过程约束的递进式训练思路，先在相对受控的文本环境中强化模型的基础演绎能力，再在图文联合输入下引入步骤级反馈，以提升模型在复杂几何问题中的推理稳定性。

\begin{figure}[H]
  \centering
  \includegraphics[width=1.0\textwidth]{figures/c4_framework.pdf}
  \caption{GeoTSR 总体框架}
  \label{fig:c4_framework}
\end{figure}

第一阶段为纯文本环境下的长链条推理增强。考虑到多模态模型在直接处理图文联合输入时，需要同时完成视觉信息理解和逻辑推导，训练目标较为复杂，本阶段首先将几何图像转化为结构化文本谓词，构建纯文本训练输入。这样可以降低视觉信息对推理学习的干扰，使模型将训练重点集中于条件整理、关系推导和步骤组织。通过轨迹级规则反馈，本阶段引导模型生成规范的思维链推理过程，增强其在长链条几何任务中的基础推理能力，并为后续多模态训练提供初始化基础。

第二阶段为多模态条件下的过程级强化。在第一阶段形成基础推理能力后，本阶段重新引入题目文本与几何图像，使模型适应真实几何任务中的图文联合输入。为缓解长链推理中的中间错误累积，本阶段利用弱符号验证器对模型生成的中间结论进行局部检验，并将检验结果转化为步骤级反馈信号。通过这种过程级约束，模型能够在生成过程中及时抑制无效推理和路径偏移，从而提升多模态长链条推理的稳定性。

% \xsection{核心组件与基础算法}{Core Components and Foundational Algorithms}

\xsubsection{弱符号验证器}{Weak Symbolic Validator}

两阶段强化学习框架中的核心组件是弱符号验证器，该验证器与传统强符号系统存在本质区别。以 AlphaGeometry~\cite{trinh2024alphageometry} 为代表的强符号求解模型通常具备较为完备的公理库，能够通过自主搜索完成几何问题的端到端全局推理；而本章所采用的弱符号验证器并不以独立求解问题为目标，也不具备完整的全局搜索能力。其主要作用在于对模型生成的单步推理结果进行局部合法性检验，并在此基础上返回相应的符号规则反馈。具体而言，弱符号验证器由解析器、谓词融合器和答案评估器三个模块构成，三者协同完成从结构校验、语义提取、谓词表示构建到结果一致性判定的处理流程，其整体结构如图~\ref{fig:c4_weak_symbolic}所示。

\begin{figure}[htbp]
  \centering
  \includegraphics[width=1.0\textwidth]{figures/c4_weak_symbolic.pdf}
  \caption{弱符号验证器整体结构}
  \label{fig:c4_weak_symbolic}
\end{figure}

为避免后文奖励符号歧义，在本节中统一约定：小写符号 $r_{\ast}$ 表示步骤级奖励，用于评价当前推理步骤 $s_t$ 的局部质量；大写符号 $R_{\ast}$ 表示轨迹级奖励，用于评价完整推理轨迹的整体表现。若某一奖励项同时具有步骤级与轨迹级两种形式，则分别以对应的小写和大写符号表示。

解析器作为系统介入的第一环节，承担着格式校验与语义转化的双重任务。首先，解析器利用正则表达式对模型输出的原始字符串进行结构扫描，检验其是否严格包裹在预设的 \texttt{<step></step>} 等规范标签内。该步骤的检验结果直接映射为格式规范奖励，步骤级记为 $r_{\mathrm{fmt}}$，轨迹级汇总记为 $R_{\mathrm{fmt}}$。在格式合法的前提下，解析器进一步利用前一章节构建的基于轻量级大语言模型的转换工具链，将提取出的自然语言推导步骤 $s_t$ 映射为结构化谓词集合步骤 $P_t$。例如，将“由于 AB 等于 AC”转化为 “Equals(LengthOf(Line(A,B)), LengthOf(Line(A,C)))” 的谓词表示。此外，解析器在后台维护实体映射表，以处理点、线、圆等实体的规范化命名，确保后续符号表示的严格一致。

谓词融合器负责维护解题过程中的动态状态空间。融合器利用字典存储当前已知的几何实体属性，并引入并查集（Disjoint Set Union, DSU）算法管理等价关系。对于几何推理中的边长相等、角度相等及平行等核心关系，融合器通过并查集维护等价类集合 $\mathcal{E} = \{E_1, E_2, \dots, E_m\}$。当新步骤 $P_t$ 产生时，融合器执行如下判定逻辑：若 $P_t$ 中的等价关系可通过当前并查集的查找操作验证，则判定为冗余；若 $P_t$ 与已有等价类冲突，则标记为逻辑冲突。该模块输出轨迹可验证性奖励 $R_{\mathrm{ver}}$、步骤可验证奖励 $r_{\mathrm{ver}}$ 以及冗余惩罚 $r_{\mathrm{red}}$。

答案评估器承担最终目标验证的职能。由于方法中已通过步骤级反馈约束了中间推理的合法性与非冗余性，答案评估器不再进行复杂的中间逻辑闭包推导，而是专注于计算当前轨迹得出的最终结论与标准答案之间的一致性。该模块输出答案匹配奖励 $R_{\mathrm{ans}}$，用以从宏观层面客观评估求解的正确性。

综上，弱符号验证器通过将几何推理过程转化为可计算的操作，为强化学习提供了低成本且客观的反馈机制。通过解析、融合与评估三个环节的协同作用，验证器能够将推理过程中的局部质量、逻辑一致性与结果正确性映射为对应的奖励信号，从而为模型的长链条推理优化提供支撑。

\xsubsection{基于组相对策略优化}{Group Relative Policy Optimization}

近年来，强化学习在大语言模型推理能力增强中的作用受到广泛关注~\cite{chen2024improving,zhang2024restmcts,liu2019deeprl}，并逐渐成为复杂推理任务中的重要优化手段。DeepSeek-R1~\cite{guo2025deepseekr1} 工作表明，基于强化学习的训练方式能够有效提升模型的长链条推理能力。受此启发，并结合大规模长文本训练对显存开销和训练稳定性的要求，本章并未采用依赖额外价值模型的近端策略优化（Proximal Policy Optimization, PPO）算法\cite{schulman2017ppo}，而是选用组相对策略优化（Group Relative Policy Optimization, GRPO）算法~\cite{guo2025deepseekr1}，对两阶段训练过程中的策略模型进行迭代更新。

对于长链条几何推理任务，同一道题往往可以生成多条不同的推理路径。例如，对于同一角度或长度的求解，模型可能从三角形性质、平行线关系或相似关系等不同角度展开推导。GRPO 通过对同一问题下多个候选解答进行组内比较来估计相对优势，无需额外训练价值模型，因而更适合几何推理中常见的同题多路径训练场景，也有利于降低长序列训练中的计算开销并保持训练过程的稳定性。

具体而言，对于每个输入问题 $q$，当前策略模型 $\pi_{\theta}$ 会并行采样 $G$ 条不同的推理轨迹 $\{\tau_1, \tau_2, \dots, \tau_G\}$，其中 $\tau_i=(a_{i,1},a_{i,2},\dots,a_{i,T_i})$ 表示第 $i$ 条推理轨迹。在计算出每条轨迹对应的绝对总奖励 $R_i$ 后，算法通过对组内奖励进行标准化来估算每条轨迹的相对优势 $A_i$：
\begin{equation}
  A_i = \frac{R_i - \text{mean}(\mathbf{R})}{\text{std}(\mathbf{R})},
\end{equation}
其中，$\mathbf{R}$ 为当前问题下 $G$ 个采样轨迹的绝对奖励集合。这里的奖励反映的是整条推理轨迹的整体质量，因此优势值衡量的是某条解题路径相对于同组其他路径的优劣程度。

在此基础上，引入 KL 散度约束以防止模型在优化中发生策略崩塌。GRPO 完整的代理优化目标定义为：
\begin{equation} \label{eq:grpo_loss}
  \begin{aligned}
    \max_{\theta}\ \mathbb{E}_{q \sim P(Q), \{\tau_i\}_{i=1}^G \sim \pi_{\theta_{\text{old}}}} \Bigg[ \frac{1}{G} \sum_{i=1}^G \Bigg( & \sum_{t=1}^{T_i} \min \Big( \rho_{i,t} A_i, \text{clip}(\rho_{i,t}, 1-\epsilon, 1+\epsilon) A_i \Big) \\
                                                                                                                                      & - \beta\, \mathbb{D}_{\mathrm{KL}}(\pi_{\theta} \Vert \pi_{\mathrm{ref}}) \Bigg) \Bigg],
  \end{aligned}
\end{equation}
其中，$\rho_{i,t} = \frac{\pi_\theta(a_{i,t}|s_{i,t})}{\pi_{\theta_{\text{old}}}(a_{i,t}|s_{i,t})}$ 为重要性采样比率，$\epsilon$ 为裁剪超参数，$\beta$ 为 KL 惩罚系数，$\pi_{\mathrm{ref}}$ 为参考策略。

在实际训练中，GRPO 通过最大化式~\eqref{eq:grpo_loss} 所示的代理目标，对当前策略模型 $\pi_{\theta}$ 进行更新，使其逐步提高高质量推理轨迹的生成概率，并抑制质量较差的候选路径。在本章方法中，GRPO 被统一用于两个阶段的强化训练：第一阶段用于提升模型在纯文本环境下的长链条推理能力，第二阶段则在此基础上继续对多模态输入条件下的策略模型进行优化。为实现阶段间衔接，第一阶段训练后得到的策略模型记为 $\pi_{\theta}^{(1)}$，并将其作为第二阶段训练的初始化模型，即令 $\pi_{\theta,0}^{(2)}=\pi_{\theta}^{(1)}$。由此，模型可以在第一阶段所形成的文本推理基础上，进一步适应图文联合输入下的几何推理任务。

% ==========================================================
\xsection{文本推理增强}{Textual Reasoning Enhancement}

第一阶段旨在先行建立模型对长链条几何推理过程的基本建模能力。为减少图像因素对训练的干扰，本阶段将几何问题统一表示为纯文本输入，使模型在不依赖视觉特征的条件下学习题目条件解析、中间结论生成以及推理步骤衔接。该阶段得到的模型可作为后续多模态训练的初始化基础。以下小节将详细介绍第一阶段的训练数据构造、奖励设计以及优化目标。

% \begin{figure}[H]
%     \centering
%     \includegraphics[width=0.80\textwidth]{figures/c4_phase1.png}
%     \caption{阶段一训练框架图}
%     \label{fig:c4_phase1}
% \end{figure}

\xsubsection{文本化推理样本构造}{Construction of Textualized Reasoning Samples}

为了确保模型在纯文本环境下能够获取完整的解题逻辑，本阶段对原始几何训练数据进行了语义重构，构建了用于强化训练的纯文本数据集。与此同时，为便于后续分析视觉输入对训练效果的影响，本文还构建了等规模的多模态对比数据集，但该部分数据不参与第一阶段强化训练。

% 纯文本训练数据集是本阶段训练的核心语料。首先包含大规模纯文本数学推理数据集 OpenR1-Math-220k \footnote{https://huggingface.co/datasets/open-r1/OpenR1-Math-220k}。该数据集是利用 DeepSeek-R1~\cite{guo2025deepseekr1}为 NuminaMath 1.5~\cite{li2024numinamath} 题库进行推理蒸馏而生成的合成语料，包含了海量问题的详细推导过程。考虑到训练效率与领域匹配度，本阶段基于题目难度与相关性对其进行了分层采样，最终筛选出约 3 万条高质量子集用作训练语料。此外，为增强几何领域专业性，本章利用前一章节构建的图文解析与形式化转换工具链，将 UniGeo~\cite{chen2022unigeo} 中筛选出的约 1 万条多模态样本统一转换为纯文本数据。视觉解析结果与自然语言条件融合成一段完整的纯文本内容作为模型输入。前者用于提供通用数学推理样本，后者用于补充几何领域的专门知识。

% 多模态对比数据集用于后续实验中评估视觉输入对训练效果的影响。该部分引入了等规模的 MathV360K~\cite{shi2024mathllava} 推理数据集，保留题目图像及其对应文本描述，作为多模态训练设置下的参照数据。该数据集旨在在样本规模相近的前提下，作为后续实验中的对照设置。

纯文本训练数据集是本阶段训练的核心语料。首先包含大规模纯文本数学推理数据集 OpenR1-Math-220k \footnote{https://huggingface.co/datasets/open-r1/OpenR1-Math-220k}。该数据集利用 DeepSeek-R1~\cite{guo2025deepseekr1} 对 NuminaMath 1.5~\cite{li2024numinamath} 题库进行推理蒸馏生成，包含大量问题的详细推导过程。考虑到训练效率与领域匹配度，本阶段基于题目难度与相关性进行分层采样，最终筛选出约 3 万条高质量子集，作为通用数学推理训练语料。此外，为增强模型对几何任务的适应性，利用前一章节构建的图文解析与形式化转换工具链，将 UniGeo~\cite{chen2022unigeo} 中筛选出的约 1 万条多模态样本统一转换为纯文本数据。视觉解析结果与自然语言条件被融合为完整的纯文本输入，用于强化几何知识与表征。

多模态对比数据集用于后续实验中评估视觉输入对训练效果的影响。该部分引入等规模的 MathV360K~\cite{shi2024mathllava} 推理数据集，保留题目图像及其对应文本描述，作为多模态训练设置下的参照数据。该数据集不参与第一阶段纯文本训练，而是在样本规模相近的前提下，用于构造后续实验中的对照设置。

\xsubsection{基于规则的轨迹级奖励}{Rule-based Trajectory Rewards}

在模态解耦后的纯文本环境下，本阶段采用公式 (\ref{eq:grpo_loss}) 的 GRPO 算法框架进行策略迭代。为约束模型输出格式并确立基础逻辑范式，本阶段设计了三个维度的轨迹级奖励信号，统一使用大写符号 $R$ 表示。轨迹级奖励在模型完成整条思维链输出后统一生成，重点评估推理轨迹的完整性、规范性与准确性。

格式规范奖励（$R_{\mathrm{fmt}}$）：该奖励要求模型遵循预设的结构化输出范式，即将整体推理过程置于 \texttt{<think></think>} 标签内，最终结论置于 \texttt{<answer></answer>} 标签中，且每一次中间推理步骤包裹在 \texttt{<step></step>} 标签之内。该约束希望在纯文本训练阶段确立细粒度的格式规范，促使模型养成步骤化输出习惯，构建稳定的推理模板。

答案奖励（$R_{\mathrm{ans}}$）：该奖励用于核查最终输出数值或表达式的准确性。通过引入数学表达式等价检验机制，系统能够识别并兼容不同但等价的数值或符号表达式，克服传统字符串精确匹配的局限性。

轨迹可验证奖励（$R_{\mathrm{ver}}$）：该奖励在模型生成完整推理轨迹后进行全局评估。系统依托解析器与谓词融合器，对轨迹内包含的所有中间步骤进行批量转换与合法性核查，并根据能够成功形式化且通过规则检验的有效步骤比例计算综合评分。引入该项奖励能够抑制脱离已知条件的逻辑幻觉和无法被解析的无效输出，保证长链条基础推演的全局合法性。

第一阶段的轨迹总奖励函数表示为：
\begin{equation}
  R^{(1)} = \lambda_1 R_{\mathrm{fmt}} + \lambda_2 R_{\mathrm{ans}} + \lambda_3 R_{\mathrm{ver}},
\end{equation}
其中，$\lambda_1, \lambda_2, \lambda_3$ 为各项奖励的权重系数。

完成本阶段训练后，收敛的模型权重将作为后续多模态强化学习的初始参数。该初始化策略的设计目的在于提前缓解长链条推导中常见的逻辑断裂问题，使后续引入视觉特征时，模型能够在已有推理能力的基础上进一步适应图文联合输入，并通过过程反馈缓解多模态长链推理中的路径偏移。

% ==========================================================
\xsection{多模态过程强化}{Multimodal Process Reinforcement}

第二阶段在恢复图文联合输入的多模态场景下，重点提升长链条几何推理的过程稳定性。本阶段聚焦图文联合输入带来的推理偏移与中间错误累积问题，通过过程级反馈约束模型的中间推理步骤，减少无效扩展，并维持模型在复杂视觉条件下的推理连续性。

\xsubsection{多模态输入}{Multimodal Inputs}

在由纯文本训练转向多模态任务时，本阶段以平面几何作为强化学习的目标领域。该领域在逻辑结构上与第一阶段保持一致，同时包含额外的视觉感知与图文理解要求。由于模型已在第一阶段形成基础推理先验，本阶段的优化重点在于引导模型将视觉信息正确融入已有推理过程，而不是重新学习完整的几何推理模式。

为构建高质量的训练环境，本阶段整合多个经典几何数据集，形成包含 2 万余条多模态几何题目的训练语料。在数据构成上，训练语料主要包括 9022 条 PGPS9K~\cite{zhang2023pgps} 数据和 7528 条 GeoQA+~\cite{cao2022geoqa_plus} 数据，并进一步加入 3002 条 Geometry3K~\cite{lu2021intergps} 数据和 1789 条 GeoSense~\cite{xu2025geosense} 数据进行扩充。通过多源数据组合，可以提升训练样本在题型、图形结构和推理路径上的多样性，为多模态强化学习提供更丰富的训练场景。

此阶段模型接收原始题目文本与几何图像的联合输入。与第一阶段主要在轨迹末端进行整体评估不同，本阶段进一步引入步骤级过程反馈：在模型生成中间推理步骤后，弱符号验证器调用解析器与谓词融合器，将当前步骤解析为结构化谓词，并与由图像解析结果和文本已知条件构成的动态条件池进行一致性比对。该机制能够在训练过程中对关键中间结论进行局部约束，从而抑制由视觉噪声或图文理解偏差引发的级联错误传播。

\xsubsection{基于规则的步骤级奖励}{Rule-based Step-level Rewards}

基于弱符号验证器的步骤级反馈，本阶段构建了多维度的步骤级奖励函数，并统一使用小写符号 $r$ 表示。与第一阶段在完整轨迹生成后进行整体评估不同，第二阶段将奖励信号细化到中间推理步骤，用于约束模型在多模态输入下的局部推导过程。具体而言，步骤级奖励主要包括以下三个维度：

格式规范奖励（$r_{\mathrm{fmt}}$）：该奖励与第一阶段的格式奖励保持一致，要求模型整体推理过程遵循 \texttt{<think>} 与 \texttt{<answer>} 标签结构，并将每个中间推理步骤置于 \texttt{<step></step>} 标签内。该约束用于保持模型在多模态环境下的步骤化输出习惯，避免视觉信息引入后造成输出结构混乱。

步骤可验证奖励（$r_{\mathrm{ver}}$）：该奖励依托解析器与谓词融合器，对当前生成的中间步骤进行局部检验。若当前步骤能够成功转化为形式化谓词，并通过动态条件池中的规则一致性检查，则系统给予正向奖励，用于鼓励模型生成可解析、可验证的中间结论。

冗余惩罚（$r_{\mathrm{red}}$）：该惩罚用于抑制循环论证和重复推导。若当前步骤得到的谓词已经在动态条件池中存在等价关系，则说明该步骤没有引入新的有效信息，系统将其判定为冗余推导并施加惩罚。

在推理过程的第 $t$ 步，步骤级中间奖励定义为：
\begin{equation}
  r_t^{(2)} = \alpha_1 r_{\mathrm{fmt}} + \alpha_2 r_{\mathrm{ver}} - \alpha_3 r_{\mathrm{red}}.
\end{equation}

当整道题目的推导轨迹 $\tau$ 结束时，系统进一步叠加最终答案匹配奖励 $R_{\mathrm{ans}}$。因此，第二阶段整条多模态推导轨迹的总奖励 $R^{(2)}(\tau)$ 表示为：
\begin{equation}
  R^{(2)}(\tau) = \sum_{t=1}^{T} r_t^{(2)} + \alpha_4 R_{\mathrm{ans}},
\end{equation}
其中，$\alpha_1$ 至 $\alpha_4$ 为各项奖励的权重系数。该轨迹总奖励随后用于 GRPO 中的组内相对优势估计，从而引导模型提高高质量多模态推理路径的生成概率。

% 上述规则奖励的作用在于为模型提供比最终答案更细粒度的训练反馈。通过同时约束输出格式、步骤合法性和推导冗余性，模型能够在生成过程中逐步减少无效中间结论，并保持较为清晰的推理路径。这有助于缓解长链条几何推理中错误难以及时纠正、并在后续步骤中持续累积的问题。

\xsubsection{优化目标与约束}{Optimization Objectives and Constraints}

第二阶段的强化训练以第一阶段得到的策略模型作为初始化模型，使模型在已有文本推理能力的基础上进一步适应图文联合输入。与第一阶段相比，本阶段的优化目标不仅关注最终答案是否正确，还通过步骤级奖励约束中间推理过程，从而缓解多模态长链推理中的路径偏移和错误累积问题。训练当中，第一阶段训练完成后的策略模型记为 $\pi_{\theta}^{(1)}$，并以此作为第二阶段训练的初始策略，即
\begin{equation}
  \pi_{\theta,0}^{(2)} = \pi_{\theta}^{(1)}.
\end{equation}

在第二阶段训练开始时，系统冻结当前初始策略的副本作为参考策略 $\pi_{\mathrm{ref}}$，用于 KL 约束；随后，当前策略 $\pi_{\theta}$ 在多模态输入和步骤级奖励的共同作用下进行更新。对于每条采样轨迹，系统首先根据式中定义的 $R^{(2)}(\tau)$ 计算轨迹总奖励，并据此完成 GRPO 的组内相对优势估计；同时，通过 KL 散度限制当前策略相对于参考策略的更新幅度，避免强化训练过程中发生过大的策略偏移。其计算形式可写为：
\begin{equation}
  \mathbb{D}_{\mathrm{KL}}(\pi_{\theta}\Vert \pi_{\mathrm{ref}})
  =
  \sum_{a}\pi_{\theta}(a\mid s)\log\frac{\pi_{\theta}(a\mid s)}{\pi_{\mathrm{ref}}(a\mid s)}.
\end{equation}

% 需要说明的是，第一阶段模型在本阶段中的作用主要是提供初始化参数，而非直接作为持续更新的训练目标。参考策略仅作为冻结分布用于约束当前策略的更新幅度。通过这种设计，模型既能够继承第一阶段形成的长链条推理能力，又能够在第二阶段充分学习图文联合输入下的推理策略，从而提升复杂几何任务中的过程稳定性。

% ==========================================================
\xsection{训练总体流程}{Overall Algorithm Pipeline}

GeoTSR 整体流程分为纯文本推理增强和多模态过程强化两个阶段，前者用于建立基础长链条推理能力，后者在此基础上进一步引入图文联合输入与步骤级反馈，算法形式的训练过程如算法 \ref{alg:two_stage} 所示。

% \begin{algorithm2e}[H]
%   % --- 解决标题格式与名称问题 ---
%   \renewcommand{\algorithmcfname}{算法} % 将 "Algorithm" 替换为 "算法"
%   \SetAlCapNameFnt{\small}              % 将 "算法 X" 标签的字体设为 small
%   \SetAlCapFnt{\small}                  % 将 caption 文本内容的字体设为 small
%   % ------------------------------

%   \setstretch{0.7} % 将行间距设置为正常值的 0.7 倍
%   \small
%   \SetAlgoLined
%   % 自定义 Input 和 Output 标签
%   \SetKwInput{KwInput}{输入}
%   \SetKwInput{KwOutput}{输出}

%   \KwInput{多模态基座模型 $\pi_{\theta}$，纯文本数据集 $\mathcal{D}_{text}$，多模态几何数据集 $\mathcal{D}_{geo}$，弱符号验证器 $\mathcal{V}$，组采样大小 $G$}
%   \KwOutput{增强后的多模态几何推理模型 $\pi_{final}$}

%   \tcp{阶段一：纯文本长链条推理增强}
%   冻结初始策略副本作为参考策略 $\pi_{\mathrm{ref}} \leftarrow \mathrm{stopgrad}(\pi_{\theta})$\;
%   \While{阶段一策略未收敛}{
%   从 $\mathcal{D}_{text}$ 中采样纯文本问题 $q$\;
%   使用采样策略 $\pi_{\theta_{\mathrm{old}}}$ 并行采样 $G$ 条推理轨迹 $\{\tau_i\}_{i=1}^G$\;
%   \For{每条轨迹 $\tau_i$}{
%   调用验证器 $\mathcal{V}$ 进行全局解析与核查，计算轨迹总奖励：
%   $R_i^{(1)} = \lambda_1 R_{\mathrm{fmt}} + \lambda_2 R_{\mathrm{ans}} + \lambda_3 R_{\mathrm{ver}}$\;
%   }
%   根据组内奖励 $\{R_i^{(1)}\}_{i=1}^G$ 计算相对优势 $A_i$\;
%   引入基于 $\pi_{\mathrm{ref}}$ 的 KL 散度约束，利用 GRPO 更新策略参数 $\theta$\;
%   }
%   得到阶段一策略模型 $\pi_{\theta}^{(1)} \leftarrow \pi_{\theta}$\;

%   \tcp{阶段二：多模态过程强化}
%   以阶段一模型初始化第二阶段策略 $\pi_{\theta} \leftarrow \pi_{\theta}^{(1)}$\;
%   冻结第二阶段初始策略副本作为参考策略 $\pi_{\mathrm{ref}} \leftarrow \mathrm{stopgrad}(\pi_{\theta})$\;
%   \While{阶段二策略未收敛}{
%   从 $\mathcal{D}_{geo}$ 中采样多模态几何问题 $q_{vis}$\;
%   使用采样策略 $\pi_{\theta_{\mathrm{old}}}$ 并行采样 $G$ 条推理轨迹 $\{\tau_j\}_{j=1}^G$\;
%   \For{每条轨迹 $\tau_j$}{
%   初始化轨迹总奖励 $R_j^{(2)} \leftarrow 0$\;
%   \For{轨迹中的每个中间推理步骤 $t=1$ \KwTo $T_j$}{
%   调用验证器 $\mathcal{V}$ 进行格式校验、步骤检验与冗余判断\;
%   计算当前步骤级奖励：
%   $r_{t,j}^{(2)} = \alpha_1 r_{\mathrm{fmt}} + \alpha_2 r_{\mathrm{ver}} - \alpha_3 r_{\mathrm{red}}$\;
%   $R_j^{(2)} \leftarrow R_j^{(2)} + r_{t,j}^{(2)}$\;
%   }
%   调用验证器 $\mathcal{V}$ 比对最终答案\;
%   叠加最终答案匹配奖励：
%   $R_j^{(2)} \leftarrow R_j^{(2)} + \alpha_4 R_{\mathrm{ans}}$\;
%   }
%   根据组内奖励 $\{R_j^{(2)}\}_{j=1}^G$ 计算相对优势 $A_j$\;
%   引入基于 $\pi_{\mathrm{ref}}$ 的 KL 散度约束，利用 GRPO 更新策略参数 $\theta$\;
%   }
%   $\pi_{final} \leftarrow \pi_{\theta}$\;
%   \caption{GeoTSR 两阶段长链条推理增强算法}
%   \label{alg:two_stage}
% \end{algorithm2e}

\begin{algorithm2e}[H]
  \renewcommand{\algorithmcfname}{算法}
  \SetAlCapNameFnt{\small}
  \SetAlCapFnt{\small}

  \setstretch{0.7}
  \small
  \SetAlgoLined

  \SetKwInput{KwInput}{输入}
  \SetKwInput{KwOutput}{输出}

  \KwInput{基座策略模型 $\pi_{\theta}$，纯文本数据集 $\mathcal{D}_{text}$，多模态几何数据集 $\mathcal{D}_{geo}$，弱符号验证器 $\mathcal{V}$，组采样大小 $G$}
  \KwOutput{增强后的多模态几何推理模型 $\pi_{final}$}

  $\pi_{\mathrm{ref}} \leftarrow \mathrm{stopgrad}(\pi_{\theta})$ \tcp*[r]{阶段一：纯文本增强}
  \While{阶段一策略未收敛}{
  从 $\mathcal{D}_{text}$ 中采样问题 $q$\;
  使用 $\pi_{\theta_{\mathrm{old}}}$ 采样 $G$ 条轨迹 $\{\tau_i\}_{i=1}^{G}$\;

  \For{每条轨迹 $\tau_i$}{
  调用 $\mathcal{V}$ 对 $\tau_i$ 进行轨迹级检验\;
  计算轨迹奖励 $R_i^{(1)}$\;
  }

  根据 $\{R_i^{(1)}\}_{i=1}^{G}$ 计算组内优势 $\{A_i\}_{i=1}^{G}$\;
  基于 $\pi_{\mathrm{ref}}$ 的 KL 约束，使用 GRPO 更新 $\theta$\;
  }

  $\pi_{\theta}^{(1)} \leftarrow \pi_{\theta}$\;

  $\pi_{\theta} \leftarrow \pi_{\theta}^{(1)},\ \pi_{\mathrm{ref}} \leftarrow \mathrm{stopgrad}(\pi_{\theta})$ \tcp*[r]{阶段二：多模态强化}
  \While{阶段二策略未收敛}{
  从 $\mathcal{D}_{geo}$ 中采样多模态问题 $q_{vis}$\;
  使用 $\pi_{\theta_{\mathrm{old}}}$ 采样 $G$ 条轨迹 $\{\tau_j\}_{j=1}^{G}$\;

  \For{每条轨迹 $\tau_j$}{
  $R_j^{(2)} \leftarrow 0$\;

  \For{每个中间步骤 $t=1$ \KwTo $T_j$}{
  调用 $\mathcal{V}$ 对步骤 $t$ 进行检验\;
  计算步骤奖励 $r_{t,j}^{(2)}$\;
  $R_j^{(2)} \leftarrow R_j^{(2)} + r_{t,j}^{(2)}$\;
  }

  调用 $\mathcal{V}$ 检验最终答案\;
  $R_j^{(2)} \leftarrow R_j^{(2)} + \alpha_4 R_{\mathrm{ans}}$\;
  }

  根据 $\{R_j^{(2)}\}_{j=1}^{G}$ 计算组内优势 $\{A_j\}_{j=1}^{G}$\;
  基于 $\pi_{\mathrm{ref}}$ 的 KL 约束，使用 GRPO 更新 $\theta$\;
  }

  $\pi_{final} \leftarrow \pi_{\theta}$\;

  \caption{GeoTSR 两阶段长链条推理增强算法}
  \label{alg:two_stage}
\end{algorithm2e}

从训练范式上看，本章方法的优势在于将原本需要在推理阶段完成的部分搜索与校验过程前移到训练阶段。传统复杂推理方法常借助 MCTS~\cite{coulom2006mcts} 或 Best-of-N~\cite{stiennon2006bestofn} 等测试期计算机制，通过在线搜索或多次采样提升最终结果的正确率，但这类方法通常会带来更高的计算开销和更长的推理时延。相比之下，本章方法利用弱符号验证器提供的奖励信号以及 GRPO 的组内比较机制，在训练过程中逐步将有效的推理模式内化到模型参数中。这样，在测试阶段模型只需进行一次前向自回归生成，即可输出完整的推理过程，而不再依赖额外的路径搜索或外部调度。因此，该训练范式在提升推理能力的同时降低了在线计算负担，更适合部署在参数规模受限的多模态模型上。

% ==========================================================
\xsection{实验与分析}{Experiments and Analysis}

本节对两阶段长链条推理增强方法进行实验分析。实验采用几何领域通用评测集 MathVerse~\cite{zhang2024mathverse} 和前文构建的长链条几何推理评测基准 GeoLStep，分别考察模型在通用数学视觉推理场景和长链条几何推理场景下的表现。对比实验选取多种主流多模态大模型作为参照，并从总体准确率、推理链长度变化以及中间过程质量等方面进行分析。随后，本节通过消融实验和案例分析，进一步讨论两阶段框架中不同模块对模型推理过程的作用。

\xsubsection{实验设置}{Experimental Setup}

\subsubsection{实验平台及配置}

实验基于高性能计算服务器展开，模型训练与推理均在 PyTorch 深度学习框架下完成。服务器硬件核心配置包括 NVIDIA A800 计算显卡，能够为 8B 参数规模的多模态模型提供充足的显存空间，支持长文本序列的强化学习训练。软件环境采用 Ubuntu 22.04 操作系统，并配置了 CUDA 12.1。具体软硬件平台及版本信息如表 \ref{tab:platform} 所示。

\begin{table}[H]
  \centering
  \small
  \caption{实验平台及配置相关信息}
  \label{tab:platform}
  \renewcommand{\arraystretch}{0.85}

  \begin{tabularx}{\textwidth}{Y Y}
    \toprule
    软硬件平台   & 版本/型号                         \\
    \midrule
    操作系统    & Ubuntu 22.04                  \\
    处理器     & Intel Xeon Gold 6248R         \\
    显卡      & NVIDIA A800 (80GB) $\times$ 4 \\
    CUDA    & 12.1                          \\
    Python  & 3.10                          \\
    PyTorch & 2.1.0                         \\
    \bottomrule
  \end{tabularx}
\end{table}

\subsubsection{参数设置}

在两阶段强化学习训练中，本章主要采用GRPO强化学习算法。针对几何推理的长链条特性，设置输入序列最大长度为 2048，以确保思维链（CoT）生成的完整性。强化学习的核心参数包括组采样大小 $G$ 设置为 16，KL 散度惩罚系数 $\beta$ 设置为 0.04。具体训练参数与各阶段奖励权重设置如表 \ref{tab:hyperparameters} 所示。

\begin{table}[htbp]
  \centering
  \small
  \caption{强化学习实验参数设置}
  \label{tab:hyperparameters}
  \renewcommand{\arraystretch}{0.85}
  \begin{tabularx}{\textwidth}{Y Y Y}
    \toprule
    实验参数                  & 参数值                & 参数说明         \\
    \midrule
    batch\_size           & 128                & 批处理样本数       \\
    learning\_rate        & $5 \times 10^{-6}$ & 初始学习率        \\
    max\_seq\_len         & 2048               & 最大输入序列长度     \\
    group\_size ($G$)     & 16                 & GRPO 组采样数量   \\
    KL\_penalty ($\beta$) & 0.04               & 策略散度惩罚系数     \\
    optimizer             & AdamW              & 损失函数优化器      \\
    $\lambda_1, \alpha_1$ & 1.0                & 格式规范奖励权重     \\
    $\lambda_2, \alpha_4$ & 2.0                & 答案匹配奖励权重     \\
    $\lambda_3, \alpha_2$ & 1.5                & 轨迹/步骤可验证奖励权重 \\
    $\alpha_3$            & 0.5                & 冗余惩罚权重       \\
    \bottomrule
  \end{tabularx}
\end{table}

\subsubsection{数据集及评估指标}

实验中两个评测基准的具体组成及对应评价指标如下文所述。MathVerse 主要用于评估模型在不同图文信息条件下的多模态数学推理能力，GeoLStep 则用于考察模型在长链条平面几何任务中的答案求解与过程推理表现。二者分别对应不同的评测维度，因此本文针对其数据特点设置相应的评价指标。

MathVerse~\cite{zhang2024mathverse} 是综合性多模态数学评测基准，包含 2,612 道涵盖平面几何、立体几何及函数等领域的高质量题目。其核心特点是将每个问题转化为六个具有不同图文信息配比的版本。在本实验中，选取其中四类具有代表性的子类作为核心评估指标，分别为：文本主导（Text-dominant, TD）、仅文本（Text-only, TO）、视觉主导（Vision-dominant, VD）以及仅视觉（Vision-only, VO）。实验最终采用上述四项指标的准确率及其算术平均值（Avg）来综合评价模型的性能表现。

GeoLStep 长步骤平面几何数据基准包含 2097 道样本，其构建过程及评估指标定义见前一章节描述。实验沿用最终答案准确率（$\mathrm{Acc}_{\text{ans}}$）与过程正确率（$\mathrm{Acc}_{\text{step}}$）作为核心评估项。前者衡量模型得出正确结论的综合性能，后者则侧重于量化模型在长链条推理路径下的推理稳定性和步骤合法性。

\subsubsection{对比方法}

为验证本章提出的两阶段增强方法在不同模型架构下的通用性与有效性，实验基于 Qwen3-VL-8B~\cite{bai2025qwen3vl} 与 InternVL-3-8B~\cite{zhu2025internvl3} 两种代表性的开源基座模型（以下称为 Base）设定了以下对比基线与消融模型：

Base-CoT：基座模型的零样本思维链基线，代表模型原始的多模态推理水平。

Base-SFT：在基座模型基础上利用 MathV360K~\cite{shi2024mathllava} 和阶段二训练所使用的等量几何推理数据集进行监督微调的模型，与用于对比传统有监督微调与强化学习训练在逻辑任务上的性能差异。

Base-Multi：两阶段均采用多模态数据进行强化的对比模型。其中第一阶段采用多模态输入，第二阶段采用几何专项数据。该设置用于对比不进行模态解耦、直接采用多模态强化训练时，模型在长链条推理任务中的表现。

Base-GeoTSR：本章提出的完整两阶段框架。该模型首先在纯文本环境下建立逻辑先验，再通过多模态几何数据进行多粒度的过程级强化。


\xsubsection{实验结果}{Experimental Results}

\subsubsection{对比实验}
在 MathVerse 基准上的实验结果如表 \ref{tab:mathverse_results} 所示。整体而言，本章提出的两阶段增强模型 Qwen3-GeoTSR 展现出较强的性能优势，以 67.8\% 的平均准确率超过了基础的零样本思维链（CoT）与传统监督微调（SFT）基线，验证了本章所提框架在复杂长链条逻辑推演中的有效性。

\begin{table}[H]
  \centering
  \small
  \caption{基于 Qwen3-VL-8B 的 MathVerse 求解准确率 (\%) 对比}
  \label{tab:mathverse_results}
  \begin{tabularx}{\textwidth}{l @{\extracolsep{\fill}} Y Y Y Y Y}
    \toprule
    模型                    & TD            & TO            & VD            & VO            & Avg           \\
    \midrule
    Qwen3-CoT             & 72.2          & 66.1          & 56.6          & 48.1          & 60.8          \\
    Qwen3-SFT             & 74.5          & 68.2          & 58.7          & 50.3          & 62.9          \\
    % Qwen3-Text & 77.5 & 71.2 & 60.8 & 51.6 & 65.8 \\
    % Qwen3-Multi & 74.1 & 67.7 & 61.2 & 53.8 & 64.2 \\
    Qwen3-Multi           & 76.3          & 68.8          & \textbf{64.2} & \textbf{55.9} & 66.3          \\
    \textbf{Qwen3-GeoTSR} & \textbf{79.6} & \textbf{74.1} & 62.9          & 54.4          & \textbf{67.8} \\
    \bottomrule
  \end{tabularx}

  \vspace{1ex}
  \raggedright
  \footnotesize 注：TD（文本主导）、TO（纯文本）、VD（视觉主导）与 VO（仅视觉）。
\end{table}

具体对比来看，相较于代表原始多模态推理能力的 Qwen3-CoT，Qwen3-GeoTSR 在所有细分子集中均实现了提升，平均准确率提高了 7.0\%。更重要的是，与基于同等规模数据构建的传统监督微调基线相比，本章方法在多个维度上取得了进一步提升。这主要得益于引入GRPO强化学习算法，并通过弱符号验证器提供较为密集的规则反馈，从而缓解传统 SFT 机械模仿模式在面对复杂逻辑任务时容易产生的学习瓶颈。特别是在逻辑推演要求较高的文本主导（TD）与纯文本（TO）任务中，Qwen3-GeoTSR 的准确率分别达到 79.6\% 与 74.1\%，相较于 Qwen3-SFT 分别提升了 5.1 与 5.9\%，说明强化学习有助于构建更稳健的推理底座。

此外，表 \ref{tab:mathverse_results} 中还展示了 Qwen3-Multi 作为未进行阶段一模态解耦的对比模型表现。虽然其在纯视觉感知指标（VD 与 VO）上略有优势，分别比 Qwen3-GeoTSR 高出 1.3 与 1.5\%，但整体平均准确率仍低于 Qwen3-GeoTSR 方法。关于纯文本逻辑解耦策略与多模态同步强化之间的比较，将在后续的消融实验小节中进行具体探讨。

表 \ref{tab:geolstep_combined}展示了 GeoLStep 数据集的分桶评测结果，
%  Base-GeoTSR 
经过两阶段增强方法训练后的模型在应对超长链条推理时展现出卓越的稳定性。以 Qwen3-VL-8B 为例，在难度最高的 13--26 步推理任务中，Qwen3-GeoTSR 的答案准确率 $\mathrm{Acc}_{\text{ans}}$ 从 SFT 基线的 20.5\% 提升至 27.8\%，过程正确率 $\mathrm{Acc}_{\text{step}}$ 更是实现了从 48.5\% 到 65.5\% 的飞跃。InternVL 基座模型在同等难度的长链条任务下，答案准确率也从 13.5\% 提升至 22.8\%。实验数据充分地证明了，相较于传统模型在长链条推理下易发生的逻辑崩溃现象，在强化学习当中通过引入步骤级过程强化能够有效抑制错误累积。即使在推理链条大幅增长的情况下，模型依然能保持较高的逻辑连贯性与解题成功率。

\begin{table}[H]
  \centering
  \small
  \caption{不同基座模型在 GeoLStep 数据集的分桶评测性能对比 (\%)}
  \label{tab:geolstep_combined}
  \begin{tabularx}{\textwidth}{l Y Y Y Y Y Y Y Y}
    \toprule
    \multirow{2}{*}{模型}      & \multicolumn{2}{c}{1 -- 4 Steps} & \multicolumn{2}{c}{5 -- 12 Steps} & \multicolumn{2}{c}{13 -- 26 Steps} & \multicolumn{2}{c}{All Steps}                                                                                                                           \\
    \cmidrule(lr){2-3} \cmidrule(lr){4-5} \cmidrule(lr){6-7} \cmidrule(lr){8-9}
                             & $\mathrm{Acc}_{\text{ans}}$      & $\mathrm{Acc}_{\text{step}}$      & $\mathrm{Acc}_{\text{ans}}$        & $\mathrm{Acc}_{\text{step}}$  & $\mathrm{Acc}_{\text{ans}}$ & $\mathrm{Acc}_{\text{step}}$ & $\mathrm{Acc}_{\text{ans}}$ & $\mathrm{Acc}_{\text{step}}$ \\
    \midrule
    % \multicolumn{9}{l}{\textit{基座模型 Qwen3-VL-8B}} \\
    Qwen3-CoT                & 77.5                             & 89.5                              & 55.2                               & 76.6                          & 16.5                        & 42.9                         & 55.9                        & 75.6                         \\
    Qwen3-SFT                & 78.2                             & 90.1                              & 56.8                               & 78.6                          & 20.5                        & 48.5                         & 57.6                        & 77.7                         \\
    % Qwen3-Multi & 79.1 & 91.2 & 58.9 & 80.5 & 24.2 & 54.2 & 59.6 & 79.9 \\
    \textbf{Qwen3-GeoTSR}    & \textbf{79.8}                    & \textbf{92.5}                     & \textbf{60.5}                      & \textbf{82.4}                 & \textbf{27.8}               & \textbf{65.5}                & \textbf{61.2}               & \textbf{82.8}                \\

    \midrule
    % \multicolumn{9}{l}{\textit{基座模型 InternVL-3-8B}} \\
    InternVL-CoT             & 72.8                             & 63.7                              & 46.3                               & 48.3                          & 10.4                        & 43.2                         & 51.4                        & 49.8                         \\
    InternVL-SFT             & 73.5                             & 65.2                              & 48.1                               & 50.4                          & 13.5                        & 46.5                         & 53.0                        & 52.0                         \\
    % InternVL-Multi & 74.8 & 68.5 & 51.2 & 55.8 & 16.2 & 51.2 & 55.6 & 56.8 \\
    \textbf{InternVL-GeoTSR} & \textbf{76.0}                    & \textbf{72.5}                     & \textbf{54.5}                      & \textbf{62.4}                 & \textbf{22.8}               & \textbf{61.5}                & \textbf{58.8}               & \textbf{63.1}                \\
    \bottomrule
  \end{tabularx}
\end{table}

\subsubsection{消融实验}

为了深入剖析本章提出的两阶段增强方法中关键设计对模型最终性能的贡献，本小节分别针对模态解耦策略、奖励机制的权重参数，以及弱符号验证器的步骤级反馈机制进行了详细的消融分析。

\paragraph{模态解耦策略消融}

为验证阶段一解耦训练和阶段二多模态规则反馈策略的必要性，本节在 MathVerse 数据集上设计了分组消融实验。具体而言，将第一阶段未进行模态解耦、直接采用多模态数据强化的 Qwen3-Multi，与采用纯文本逻辑增强的 Qwen3-GeoTSR 进行系统对比，并同时列出二者仅经历第一阶段训练、去除第二阶段几何专项强化，即 w/o Geo 设置下的性能表现。具体的消融对比结果如表 \ref{tab:ablation_modality} 所示。

\begin{table}[H]
  \centering
  \small
  \caption{不同模态强化策略在 MathVerse 数据集上的准确率 (\%) 消融对比}
  \label{tab:ablation_modality}
  \begin{tabularx}{\textwidth}{l @{\extracolsep{\fill}} Y Y Y Y Y}
    \toprule
    模型                    & TD            & TO            & VD            & VO            & Avg                      \\
    \midrule
    Qwen3-Multi           & 75.9          & 68.2          & \textbf{64.0} & \textbf{55.4} & $65.8_{~\downarrow 2.0}$ \\
    \quad w/o Geo         & 73.7          & 67.1          & 61.0          & 53.3          & $63.7_{~\downarrow 4.1}$ \\
    \midrule
    \textbf{Qwen3-GeoTSR} & \textbf{79.6} & \textbf{74.1} & 62.9          & 54.4          & \textbf{67.8}            \\
    \quad w/o Geo         & 76.8          & 71.1          & 60.5          & 51.4          & $65.0_{~\downarrow 2.8}$ \\
    \bottomrule
  \end{tabularx}

  \vspace{1ex}
  \raggedright
  \footnotesize 注：TD（文本主导）、TO（纯文本）、VD（视觉主导）与 VO（仅视觉）。
\end{table}

结合表 \ref{tab:ablation_modality} 中的实验数据可以观察到，在仅进行第一阶段强化，即 w/o Geo 设置下，直接引入图文联合输入的策略能够较早地进行跨模态特征对齐，因此在视觉主导（VD）与仅视觉（VO）子集上的表现优于纯文本基座模型。然而，这种未解耦的同步训练会在一定程度上分散模型对复杂逻辑链条的学习能力，使其在逻辑要求更高的文本主导（TD）任务中落后于纯文本基座（73.7\% vs 76.8\%）。

经过完整的两阶段训练后，两种策略的差异进一步显现。未解耦的 Qwen3-Multi 虽然在仅视觉（VO）任务上取得了 55.4\% 的最高准确率，但在高度依赖逻辑推演的文本主导（TD）任务中仍存在局限，准确率为 75.9\%。相比之下，采用纯文本逻辑增强作为第一阶段训练基础的 Qwen3-GeoTSR 在逻辑推演上表现更优，其 TD 准确率达到 79.6\%，并最终以 67.8\% 的平均准确率取得最优结果。这表明，第一阶段的纯文本解耦虽然暂时没有利用视觉信息，但能够减少视觉特征对逻辑学习过程的干扰，有助于模型形成更稳定的推理能力。不同设置下平均准确率下降幅度的对比结果进一步说明，先增强逻辑推理能力、再进行跨模态对齐与几何强化的训练策略具有必要性。

\paragraph{奖励权重参数分析}

在第二阶段的多模态过程强化中，步骤级奖励权重的设置会影响弱符号验证器反馈信号的作用方式，并进一步影响模型生成推理轨迹的稳定性。为分析所设权重的合理性，本节在 GeoLStep 评测基准的 13--26 步长链条子集上，考察两个关键奖励参数对模型性能的影响，如图 \ref{fig:c4_weight_ablation} 所示。


\begin{figure}[htbp]
  \centering
  \begin{minipage}[t]{0.48\textwidth}
    \centering
    \includegraphics[width=\textwidth]{figures/c4_weight_sensitivity_analysis_a.png}
    \vspace{0.1cm}
    \centerline{\small (a) 冗余惩罚权重 $\alpha_3$ 的影响}
  \end{minipage}
  \hfill
  \begin{minipage}[t]{0.48\textwidth}
    \centering
    \includegraphics[width=\textwidth]{figures/c4_weight_sensitivity_analysis_b.png}
    \vspace{0.1cm}
    \centerline{\small (b) 结果奖励权重 $\alpha_4$ 的影响}
  \end{minipage}
  \caption{第二阶段奖励权重参数对性能的影响}
  \label{fig:c4_weight_ablation}
\end{figure}

图 \ref{fig:c4_weight_ablation} (a) 展示了冗余惩罚权重 $\alpha_3$ 对长链条推理稳定性的影响。当不施加冗余惩罚，即 $\alpha_3=0.0$ 时，模型更容易生成重复性步骤或信息量较低的推导，导致真正推动解题进程的步骤比例下降，因此答案准确率与过程正确率均处于相对较低水平。随着 $\alpha_3$ 逐渐提升至 0.5 附近，重复和无效步骤得到一定抑制，模型能够在继续探索解题路径的同时减少无意义的展开，整体性能达到较优水平。然而，当惩罚权重继续增大并接近 1.0 时，模型会变得过于保守，一些必要的中间推导也可能被抑制，导致推理链条无法充分展开。在这种情况下，答案准确率下降更为明显；而由于模型保留下来的少量步骤通常仍较为可靠，过程正确率的下降幅度相对较小。上述结果说明，冗余惩罚有助于减少长链条推理中的无效推导，但其权重需要控制在适当范围内。

图 \ref{fig:c4_weight_ablation} (b) 进一步展示了结果奖励权重 $\alpha_4$ 的影响。随着 $\alpha_4$ 的提升，答案准确率与过程正确率整体呈现先升后降的趋势。当 $\alpha_4$ 设定在 2.0 左右时，模型既能关注最终答案是否正确，也能保持对中间推理步骤的约束，综合性能达到较优水平。若继续增大该权重，训练过程会更加偏向最终答案，模型对中间推导质量的关注相对减弱。此时，模型可能在推理后期产生跳跃式推导，或在依据不足的情况下直接靠近目标答案，从而影响完整推理过程的连贯性，并导致过程正确率与答案准确率同步下降。该结果表明，在长链条几何推理任务中，最终答案奖励能够提供必要的目标引导，但仍需要与步骤级检查配合使用，才能更好地维持推理过程的可靠性。

\subsubsection{案例分析}

为了更直观地分析本章方法在长链条多模态推理中的表现，本节从 GeoLStep 数据集中选取两道典型几何题进行案例分析，分别展示模型求解成功与失败的推理过程。前者用于说明过程反馈对减少错误累积的作用，后者则揭示当前方法在处理隐性空间关系时仍存在的不足。

图 \ref{fig:c4_casestudy} 展示了一道参考推导长达 21 步的复杂几何题。如图中中间栏所示，基线模型 Qwen3-CoT 在推导初期的 Step 4 受到视觉外观误导，将图中形似直角的 $\triangle BDC$ 误判为直角三角形，并据此引入正切关系。由于该判断缺少可靠条件支撑，后续推理很快沿着错误路径展开，导致整体推理轨迹发生偏移。相比之下，本章方法GeoTSR在长链条推导中表现出更稳定的步骤推进能力。借助训练阶段引入的弱符号验证反馈，模型在生成中间结论时更倾向于依据已知条件进行推导，例如围绕倍角公式和相似三角形判定逐步展开，而不是直接依赖图形外观作出判断。最终，模型完成了 22 步推导并得到正确结果。该案例表明，过程级反馈有助于减少早期视觉误判对后续推理的影响，从而缓解长链条推理中的错误累积问题。

\begin{figure}[H]
  \centering
  \includegraphics[width=0.7\textwidth]{figures/c4_casestudy.pdf}
  \caption{长链条推理中基线方法与增强方法的过程对比}
  \label{fig:c4_casestudy}
\end{figure}

图 \ref{fig:c4_error_case} 展示了一个求解失败案例，该题的推导链条相对较短，但模型在 Step 8 中直接得出了 $AB = AF + FB = 6 + 18 = 24$ 的结论。从代数形式上看，该步骤符合线段相加关系；但这一关系成立的前提是 $A$、$F$、$B$ 三点共线且 $F$ 位于 $A$ 与 $B$ 之间。模型在生成该步骤时未能结合图像充分核对三点的实际排列顺序，因此将不满足条件的空间关系错误地转化为了线段加法。当前的弱符号验证器主要检验显式条件、等价关系与代数计算的一致性，对点序关系、线段包含关系等隐性拓扑约束覆盖不足，因此未能及时识别这一前提错误。随着该错误结论被继续用于后续正弦定理计算，最终导致求解失败。该案例说明，尽管本章方法能够在一定程度上提升长链条推理稳定性，但对于依赖图形细节和隐性空间位置的几何问题，仍需要更精细的拓扑关系验证机制作为补充。

\begin{figure}[H]
  \centering
  \includegraphics[width=0.7\textwidth]{figures/c4_error_case.pdf}
  \caption{求解失败案例}
  \label{fig:c4_error_case}
\end{figure}


% ==========================================================
\xsection{本章小结}{Summary}

本章提出了一种基于规则反馈的两阶段增强方法GeoTSR，旨在解决长链条多模态几何推理中的级联误差和推理不稳定问题。该方法通过先增强推理逻辑、后进行模态对齐的解耦策略，首先在纯文本环境下应用组相对策略优化（GRPO）提升模型的基础推理能力；随后，在多模态场景中引入弱符号验证器，提供细粒度的过程反馈，实时引导推理步骤。实验结果表明，相较于传统的基线方法，本章提出的方法在 MathVerse 和 GeoLStep 基准上均取得了优异的表现，尤其在13步以上的长链条任务中，显著提升了模型的推理稳定性和过程正确率。尽管方法在处理隐性空间拓扑关系方面存在一定的局限性，但本章验证了规则引导的强化学习在复杂几何推理任务中的有效性，为多模态神经符号系统的应用提供了有益的参考。